\section{РАЗРАБОТКА ИМИТАТОРА ПОЛЬЗОВАТЕЛЬСКИХ ЗАЯВОК}

\subsection{Анализ исходного набора данных}

Исходные данные для построения статистической модели пользовательских
заявок представлены в виде набора учётных записей SLURM
\texttt{acct\_0921-0923}, содержащего сведения о реально выполненных
задачах действующего вычислительного кластера. В соответствии с
вариантом задания из общего набора данных были отобраны записи,
удовлетворяющие условию: \texttt{UID = 50430} и
\texttt{JobName = "run\_bwa.sh"}.

% TODO: подставить реальное количество отобранных записей
В результате фильтрации было получено $N_{\text{sample}} = $ \texttt{<N>}
записей, использованных для дальнейшего статистического анализа.

\subsection{Предварительный анализ и масштабирование данных}

Анализ отобранного подмножества показал, что фактическое время
выполнения задач (\texttt{ElapsedRaw}) принимает значения порядка
% TODO: подставить реальный диапазон (часы/дни)
\texttt{<порядок величин в выборке>}, что существенно затрудняет
непосредственное воспроизведение данных интервалов времени в условиях
учебного стенда. В связи с этим было принято решение о масштабировании
исходных временных характеристик с понижающим коэффициентом
% TODO: подставить выбранный коэффициент масштабирования и обоснование
$k = $ \texttt{<k>}, что позволило перейти от исходных величин,
измеряемых часами и днями, к величинам порядка десятков-сотен секунд,
пригодным для практической демонстрации работы имитатора и диспетчера.

\subsection{Выбор модели ошибки оценки времени выполнения}

Для характеристики отклонения пользовательской оценки времени
выполнения от фактического времени выполнения задачи использовались
две альтернативные модели:
\begin{itemize}
\item абсолютное отклонение $\Delta T = T_{\text{заявки}} - T_{\text{факт}}$;
\item относительное отклонение $\tau = T_{\text{заявки}} / T_{\text{факт}}$.
\end{itemize}

% TODO: подставить обоснование выбора на основе фактических значений deltaT / Tфакт по варианту
По результатам анализа отобранной выборки было установлено, что
величина $\Delta T$ \texttt{<сопоставима с / превышает>} величину
$T_{\text{факт}}$, в связи с чем для дальнейшей работы была выбрана
модель относительного отклонения~$\tau$.

\subsection{Вычисление статистических параметров}

Для отобранной и масштабированной выборки были вычислены границы
диапазона фактического времени выполнения, а также параметры
логарифма выбранной величины ошибки оценки. Результаты вычислений
приведены в таблице~\ref{tab:stats}.

% TODO: подставить реально вычисленные значения
\begin{table}[H]
\caption{Статистические параметры выборки}
\label{tab:stats}
\centering
\begin{tabular}{|l|c|}
\hline
Параметр & Значение \\
\hline
$T_{\text{факт, min}}$, с & \texttt{<min>} \\
$T_{\text{факт, max}}$, с & \texttt{<max>} \\
$\mu$ ($\ln \tau$) & \texttt{<mu>} \\
$\sigma$ ($\ln \tau$) & \texttt{<sigma>} \\
\hline
\end{tabular}
\end{table}

\subsection{Алгоритм генерации пользовательской заявки}

Генерация каждой синтетической заявки выполняется в два этапа.

На первом этапе фактическое время выполнения задачи $T_{\text{факт}}$
моделируется как случайная величина, подчинённая нормальному
распределению, ограниченному диапазоном $[T_{\text{min}},\,
T_{\text{max}}]$, полученным в результате анализа исходных данных
(таблица~\ref{tab:stats}). Значение $T_{\text{факт}}$ используется в
качестве параметра команды \texttt{sleep} в теле формируемого
пакетного сценария SLURM.

На втором этапе пользовательская оценка времени выполнения
$T_{\text{заявки}}$ вычисляется как:
\[
T_{\text{заявки}} = T_{\text{факт}} \cdot e^{X}, \qquad X \sim
\mathcal{N}(\mu, \sigma^2),
\]
то есть отклонение оценки от фактического времени моделируется
логнормальным распределением с параметрами $\mu$ и $\sigma$,
вычисленными на этапе статистического анализа. Полученное значение
$T_{\text{заявки}}$ используется как значение параметра
\texttt{--time} пакетного сценария.

Количество узлов, запрашиваемых заявкой, выбирается случайным образом
из множества $\{1, 2, 3, 4\}$ с использованием
% TODO: подставить реально использованное распределение (равномерное / нормальное / Коши)
\texttt{<распределение>}, независимо от количества узлов, фигурирующего
в исходных данных по варианту.

\subsection{Реализация программы-имитатора}

Имитатор реализован на языке
% TODO: подставить язык реализации (Python / Bash / Go и т.д.)
\texttt{<язык реализации>}
и принимает в качестве входного параметра количество заявок $N$
(целое положительное число, не превышающее 999). Для каждой из $N$
заявок программа:
\begin{enumerate}
\item генерирует значения $T_{\text{факт}}$, $T_{\text{заявки}}$ и
количество узлов согласно описанным выше распределениям;
\item формирует пакетный сценарий SLURM с уникальным именем задачи,
указанием количества узлов (\texttt{--nodes}), предельного времени
выполнения (\texttt{--time}) и командой \texttt{sleep
<$T_{\text{факт}}$>} в качестве имитации полезной нагрузки;
\item автоматически ставит сформированную заявку в очередь диспетчера
посредством команды \texttt{sbatch}.
\end{enumerate}

\begin{lstlisting}[caption={Шаблон пакетного сценария, формируемого имитатором}]
#!/bin/bash
#SBATCH --job-name=sim_job_%j
#SBATCH --nodes=<N_nodes>
#SBATCH --time=00:00:<T_request>

sleep <T_actual>
\end{lstlisting}

\subsection{Сбор статистики поставленных в очередь заявок}

Для последующего сравнения плановых и фактических характеристик
выполненных заявок сбор статистики осуществлялся средствами SLURM с
использованием команды \texttt{sacct}, позволяющей получить сведения
о времени постановки заявки в очередь (\texttt{Submit}), запрошенном
времени выполнения (\texttt{TimelimitRaw}), фактическом времени начала
и завершения выполнения (\texttt{Start}, \texttt{End}) и фактической
продолжительности выполнения (\texttt{ElapsedRaw}). Результат
выполнения команды сохраняется в текстовый файл для последующего
анализа.

\begin{lstlisting}[caption={Сохранение статистики выполненных заявок в файл}]
sacct --starttime=<начало периода> \
      --format=JobID,JobName,Submit,Start,End,ElapsedRaw,TimelimitRaw,NNodes,State \
      --parsable2 --noheader > sim_results.csv
\end{lstlisting}